\section{The excluded excess and the full existence theorem}\label{the-excluded-excess-and-the-full-existence-theorem}

The rank-three product bound gives
\(2\cdot6\cdot2^s\leq3\cdot2^n\), so \(s\leq n-2\); nondegeneracy
gives \(s\geq2\). If \(s=n-2\), the standard tight construction for two
even adjacent orders follows from UA.SK applied to the rank-three
Boolean group. It has exactly \(3\cdot2^n\) elements. Sections 5 and 6
supply every \(s\leq n-4\). It remains to exclude \(s=n-3\).

We retain Lucchini's cyclic-core theorem {[}L98, Proposition, p.241{]}: for a
proper cyclic subgroup \(A\) of index \(t\) in a finite group \(G\),
\([G:\operatorname{core}_G(A)]\leq t^2-t\). Dividing by \(t\) gives
\([A:\operatorname{core}_G(A)]<t=[G:A]\). The complete original paper was
read and is bound as a retained primary; its proof is not rederived here. In a hypothetical group of
order \(3\cdot2^n\) and type \(\{6,2^{n-3}\}\), the second rotation
subgroup \(A=\langle bc\rangle\) has index twenty-four. For \(n\geq8\),
its order is at least thirty-two, so its core is nontrivial. Its unique
subgroup \(N\) of order two is characteristic in the cyclic core and
therefore normal in \(G\).

The quotient \(G/N\) preserves the first dihedral subgroup injectively:
that subgroup meets \(\langle b,c\rangle\) in \(\langle b\rangle\),
whereas \(N\) lies in the rotation subgroup. Because
\(N\leq\langle b,c\rangle\), the adjacent intersection in the quotient
is still exactly the image of \(\langle b\rangle\). To see this
directly, if an element of the first subgroup maps into the second
one, it already lies in that second subgroup, since its full preimage
is unchanged by adjoining \(N\). Thus the quotient is a string C-group
of order \(3\cdot2^{n-1}\) and type \(\{6,2^{n-4}\}\). Descent reaches
\(n=7\).

For that base, an independent exhaustive calculation uses the retained
complete SmallGroups catalogue of groups of order 384. For every group
\(G\) it first tests whether its exponent is divisible by sixteen. For
each remaining group, it takes one representative \(c\) from every
conjugacy class of involutions. It enumerates every involution \(b\) with
\(|bc|=16\), modulo the conjugation action of \(C_G(c)\). For each such
pair it takes every involution \(a\in C_G(c)\) with \(|ab|=6\), and tests
whether \(\langle a,b,c\rangle=G\).

Every possible marking is included: simultaneous conjugation can first
place \(c\) at its representative and then move \(b\) within its centralizer
orbit; the exhaustive centralizer search still contains the corresponding
\(a\). The complete computation visits 20,169 catalogue groups, including
1,178 whose exponent passes the first test, and 4,436 reduced adjacent
pairs. There is no generating triple, even before imposing IP. Thus
no required C-group exists at order 384. This is a catalogue-based
computer-assisted proof, not an unexplained statement that a software
package returned no example. The algorithm and counts are replayable.
The induction proves the exclusion for all \(n\geq7\), completing
Theorem 3 without a residual list of unspecified small cases.
